%
% pol-innen.tex -- Graph der Funktion f(x) = 1/(x-1) auf [0,3]
%
% (c) 2026 Damien Flury, OST Ostschweizer Fachhochschule
%
\documentclass[tikz]{standalone}
\usepackage{amsmath}
\usepackage{times}
\usepackage{txfonts}
\usepackage{pgfplots}
\usetikzlibrary{arrows,math}
\definecolor{darkblue}{rgb}{0,0,0.7}
\begin{document}
\def\skala{1}
\begin{tikzpicture}[>=latex,thick,scale=\skala]

% Achsen
\draw[->] (-0.5,0) -- (3.7,0) coordinate[label={$x$}];
\draw[->] (0,-3.3) -- (0,3.5) coordinate[label={right:$y$}];

% Achsenbeschriftung
\foreach \x in {1,2,3}{
	\draw (\x,-0.05) -- (\x,0.05);
	\node at (\x,0) [below] {$\x$};
}
\node at (0,0) [below left] {$0$};
\foreach \y in {-3,-2,-1,1,2,3}{
	\draw (-0.05,\y) -- (0.05,\y);
	\node at (0,\y) [left] {$\y$};
}

% Polstelle bei x = 1
\draw[dashed,gray] (1,-3.3) -- (1,3.5);

% Flächen unter dem Integranden, symmetrisch um die Polstelle bei x = 1
\fill[darkblue!30,opacity=0.5]
	(0,0) -- plot[domain=0:0.7,samples=100] ({\x},{1/(\x-1)}) -- (0.7,0) -- cycle;
\fill[darkblue!30,opacity=0.5]
	(1.3,0) -- plot[domain=1.3:3,samples=100] ({\x},{1/(\x-1)}) -- (3,0) -- cycle;

% linker Ast
\draw[color=darkblue,line width=1.4pt]
	plot[domain=0:0.697,samples=200] ({\x},{1/(\x-1)});
% rechter Ast
\draw[color=darkblue,line width=1.4pt]
	plot[domain=1.303:3,samples=200] ({\x},{1/(\x-1)});

\node[color=darkblue] at (2.3,1.1) {$f(x)=\dfrac{1}{x-1}$};

\end{tikzpicture}
\end{document}
